\FigureLayoutDeclare{img/2024/new_gaokao_paper_1/q11_fig1.png}{6.5cm}{32bp 84bp 74bp 28bp}
\FigureLayoutDeclare{img/2024/new_gaokao_paper_1/q12_fig1.png}{6.5cm}{0bp 0bp 0bp 0bp}
\FigureLayoutDeclare{img/2024/new_gaokao_paper_1/q17_solution_fig1.png}{8.0cm}{0bp 0bp 0bp 0bp}

\chapter{2024年新高考I卷}

\section{单选题}

\begin{problem}
已知 \(A=\{x\mid -\sqrt[3]{5}<x<\sqrt[3]{5}\}\)，\(B=\{-3,-1,0,2,3\}\)，则 \(A\cap B=\)（\quad）

\choices
    {\(\{-1,0\}\)}
    {\(\{2,3\}\)}
    {\(\{-3,-1,0\}\)}
    {\(\{-1,0,2\}\)}
\end{problem}

\begin{problem}
若 \(\dfrac{z}{z-1}=1+\i\)，则 \(z=\)（\quad）

\choices
    {\(-1-\i\)}
    {\(-1+\i\)}
    {\(1-\i\)}
    {\(1+\i\)}
\end{problem}

\begin{problem}
已知向量 \(\bs a=(0,1)\)，\(\bs b=(2,x)\)，若 \(\bs b\perp (\bs b-4\bs a)\)，则 \(x=\)（\quad）

\choices
    {\(-2\)}
    {\(-1\)}
    {\(1\)}
    {\(2\)}
\end{problem}

\begin{problem}
已知 \(\cos(\alpha+\beta)=m\)，\(\tan\alpha\tan\beta=2\)，则 \(\cos(\alpha-\beta)=\)（\quad）

\choices
    {\(-3m\)}
    {\(-\frac{m}{3}\)}
    {\(\frac{m}{3}\)}
    {\(3m\)}
\end{problem}

\begin{problem}
已知圆柱和圆锥的底面半径相等，侧面积相等，且它们的高均为 \(\sqrt3\)，则圆锥的体积为（\quad）

\choices
    {\(2\sqrt3\pi\)}
    {\(3\sqrt3\pi\)}
    {\(6\sqrt3\pi\)}
    {\(9\sqrt3\pi\)}
\end{problem}

\begin{problem}
已知函数
\(f(x)=-x^2-2ax-a\)（\(x<0\)），\(f(x)=\e^x+\ln(x+1)\)（\(x\ge 0\)）
在 \(\R\) 上单调递增，则 \(a\) 取值范围为（\quad）

\choices
    {\((-\infty,0]\)}
    {\([-1,0]\)}
    {\([-1,1]\)}
    {\([0,+\infty)\)}
\end{problem}

\begin{problem}
当 \(x\in[0,2\pi]\) 时，曲线 \(y=\sin x\) 与 \(y=2\sin\!\left(3x-\frac{\pi}{6}\right)\) 的交点个数为（\quad）

\choices
    {3}
    {4}
    {6}
    {8}
\end{problem}

\begin{problem}
已知函数 \(f(x)\) 的定义域为 \(\R\)，且
\(f(x)>f(x-1)+f(x-2)\)，
当 \(x<3\) 时 \(f(x)=x\). 则下列结论中一定正确的是（\quad）

\choices
    {\(f(10)>100\)}
    {\(f(20)>1000\)}
    {\(f(10)<1000\)}
    {\(f(20)<10000\)}
\end{problem}

\section{多选题}

\begin{problem}
为了解推动出口后的亩收入（单位：万元）情况，从该种植区抽取样本，得到推动出口后亩收入样本均值 \(x=2.1\)，样本方差 \(s^2=0.01\). 已知该种植区以往亩收入 \(X\sim N(1.8,0.1)\)，假设推动出口后的亩收入 \(Y\sim N(x,s)\)，则（\quad）（若 \(Z\sim N(u,\sigma)\)，则 \(P(Z<u+\sigma)\approx0.8413\)）

\choices
    {\(P(X>2)>0.2\)}
    {\(P(X>2)<0.5\)}
    {\(P(Y>2)>0.5\)}
    {\(P(Y>2)<0.8\)}
\end{problem}

\begin{problem}
设函数
\(f(x)=(x-1)^2(x-4)\)
，则（\quad）

\choices
    {\(x=3\) 是 \(f(x)\) 的极小值点}
    {当 \(0<x<1\) 时，\(f(x)<f(x^2)\)}
    {当 \(1<x<2\) 时，\(-4<f(2x-1)<0\)}
    {当 \(-1<x<0\) 时，\(f(2-x)>f(x)\)}
\end{problem}

\begin{problem}
造型可做成美丽的丝带，将其看作图中曲线 \(C\) 的一部分. 已知 \(C\) 过坐标原点 \(O\)，且 \(C\) 上点满足横坐标大于 \(-2\)，到点 \(F(2,0)\) 的距离与到定直线 \(x=a\)（\(a<0\)） 的距离之积为 \(4\)，则（\quad）

\bitmapfigure[max width=6.0cm]{img/2024/new_gaokao_paper_1/q11_fig1.png}

\choices
    {\(a=-2\)}
    {点 \((2\sqrt2,0)\) 在 \(C\) 上}
    {在第一象限 \(C\) 的点的纵坐标最大值为 \(1\)}
    {若点 \((x_0,y_0)\) 在 \(C\) 上，则 \(y_0\le \dfrac{4}{x_0+2}\)}
\end{problem}

\section{填空题}

\begin{problem}
设双曲线
\(\frac{x^2}{a^2}-\frac{y^2}{b^2}=1\)（\(a>0\)，\(b>0\)）
的左右焦点分别为 \(F_1\)，\(F_2\). 过 \(F_2\) 作与 \(y\) 轴平行的直线交该双曲线于 \(A\)，\(B\) 两点. 若
\(|AF_1|=13\)，\(|AB|=10\)，
则该双曲线的离心率为（\quad）.

\bitmapfigure[max width=6.0cm]{img/2024/new_gaokao_paper_1/q12_fig1.png}
\end{problem}

\begin{problem}
若曲线 \(y=\e^x+x\) 与曲线 \(y=\ln(x+1)+a\) 在点 \((0,1)\) 处有切线重合，则 \(a=\)\fillinblank{}.
\end{problem}

\begin{problem}
甲，乙两人各有四张卡片，甲为 \(\{1,3,5,7\}\)，乙为 \(\{2,4,6,8\}\). 每轮各出一张牌比较大小，先大者得 \(1\) 分，先小者得 \(0\) 分，弃置该轮牌. 四轮后甲总得分不少于 \(2\) 的概率是\fillinblank{}.
\end{problem}

\section{解答题}

\begin{problem}
记 \(\triangle ABC\) 的内角 \(A\)，\(B\)，\(C\) 对边分别为 \(a\)，\(b\)，\(c\). 已知
\(\sin C =2\cos B\)，\(a^2+b^2-c^2=\sqrt2\,ab\).
\begin{enumerate}
\item 求 \(B\)；
\item 若 \(\triangle ABC\) 的面积为 \(3+\sqrt3\)，求 \(c\).
\end{enumerate}
\end{problem}

\begin{problem}
已知 \(A(0,3)\)，\(P\left(3,\frac32\right)\) 都在椭圆
\(C:\frac{x^2}{a^2}+\frac{y^2}{b^2}=1\)（\(a>b>0\)）
上.
\begin{enumerate}
\item 求 \(C\) 的离心率；
\item 过 \(P\) 的直线 \(l\) 交 \(C\) 于另一点 \(B\)，且 \(\triangle ABP\) 的面积为 \(9\)，求 \(l\) 的方程.
\end{enumerate}
\end{problem}

\begin{problem}
如图，四棱锥 \(P-ABCD\) 中，\(PA\perp\) 底面 \(ABCD\)，且 \(PA=AC=2\)，\(BC=1\)，\(AB=\sqrt3\).

\begin{enumerate}
\item 若 \(AD\perp PB\)，证明 \(AD\parallel\) 平面 \(PBC\)；
\item 若 \(AD\perp DC\)，且二面角 \(A-CP-D\) 的正弦值为 \(\frac{\sqrt{42}}{7}\)，求 \(AD\).
\end{enumerate}

\bitmapfigure[max width=6.0cm]{img/2024/new_gaokao_paper_1/q17_fig1.png}
\end{problem}

\begin{problem}
已知函数
\(f(x)=\ln\frac{x}{2-x}+ax+b(x-1)^3\)（\(x\in(0,2)\)）.
\begin{enumerate}
\item 若 \(b=0\) 且 \(f'(x)\ge0\)，求 \(a\) 的最小值；
\item 证明：曲线 \(y=f(x)\) 关于点 \((1,a)\) 中心对称；
\item 若 \(f(x)>-2\) 当且仅当 \(1<x<2\)，求 \(b\) 的取值范围.
\end{enumerate}
\end{problem}

\begin{problem}
设 \(m\) 为正整数，数列 \(a_1\)，\(a_2\)，\(\dots\)，\(a_{4m+2}\) 为公差不为 \(0\) 的等差数列，若删去两项 \(a_i\)，\(a_j\)（\(i<j\)） 后剩余 \(4m\) 项可平均分为 \(m\) 组且每组 4 个数都成等差，则称该数列为 \((i,j)\)-可分数列.
\begin{enumerate}
\item 写出所有 \((i,j)\)，\(1\le i<j\le6\)，使 \(a_1\)，\(\dots\)，\(a_6\) 是 \((i,j)\)-可分数列；
\item 当 \(m\ge3\) 时，证明 \(a_1\)，\(\dots\)，\(a_{4m+2}\) 是 \((2,13)\)-可分数列；
\item 从 \(1\)，\(2\)，\(\dots\)，\(4m+2\) 中不放回任取 \(i<j\)，设 \(a_1\)，\(\dots\)，\(a_{4m+2}\) 是 \((i,j)\)-可分数列的概率为 \(P_m\)，证明 \(P_m>\frac18\).
\end{enumerate}
\end{problem}

